Da sowohl der \gls{tcp}-Proxy als auch der Server-Simulator Fehler zufällig erzeugen, reicht ein einzelner Durchlauf nicht aus, die Delivery-Garantien zu bestätigen, da ein fehlerfreier Durchlauf auch reines Glück bei den Zufallswerten sein könnte. Da die Hilfsservices alle aufgetretenen Ereignisse speichern, lässt sich genau zählen, wie viele Störungen während eines Laufs aufgetreten sind. Zur Auswertung werden folgende Messwerte definiert:

\begin{itemize}
    \item $N_{\text{msg}}$: Anzahl der getrennten Alarmmeldungen.
    \item $N_{\text{drop}}$: Anzahl der getrennten \gls{tcp}-Verbindungen zwischen \gls{alptwo} und dem Broker.
    \item $N_{\text{nack}}$: Anzahl der negativen Bestätigungen durch den Server-Simulator.
    \item $N_{\text{loss}}$: Anzahl der verlorene Alarmmeldungen.
    \item $N_{\text{dup}}$: Anzahl der doppelten Alarmmeldungen.
    \item $N_{\text{ooo}}$: Anzahl der Alarmmeldungen in falscher Reihenfolge.
\end{itemize}

Um zu belegen, dass die Simulation gültig ist und die Garantien eingehalten werden, wird folgendes Vorgehen gewählt:

\begin{enumerate}
    \item \textbf{Mehrere Testläufe durchführen}: Die Simulation wird mehrfach mit unterschiedlichen Zufallswerten wiederholt. Dadurch treten Verbindungsabbrüche und Serverfehler zu immer wieder anderen Zeitpunkten auf.
    \item \textbf{Fehler überprüfen}: Ein Testlauf gilt nur dann als gültig, wenn der \gls{tcp}-Proxy und der Server-Simulator Störungen verursacht haben:
          \begin{equation}
              N_{\text{drop}} > 0 \quad \text{und} \quad N_{\text{nack}} > 0
          \end{equation}
    \item \textbf{Ergebnisse überprüfen}: Nach jedem Testlauf vergleicht der \textit{Validator} die Dateien \texttt{device-results.json} und \texttt{server-results.json}, sodass gilt:
          \begin{equation}
              N_{\text{loss}} = 0, \quad N_{\text{dup}} = 0 \quad \text{und} \quad N_{\text{ooo}} = 0
          \end{equation}
    \item \textbf{Statistische Absicherung bei null Fehlern}: Da bei null aufgetretenen Fehlern keine herkömmliche Standardabweichung berechnet werden kann, wird die sogenannte \textit{Rule of Three}~\citep{hanley1983} verwendet. Wenn bei insgesamt $N$ gesendeten Alarmmeldungen kein einziger Fehler auftritt, ist die verbleibende Fehlerwahrscheinlichkeit $p$ mit einer statistischen Sicherheit von $95\,\%$ kleiner als:
          \begin{equation}
              p < \frac{3}{N}
          \end{equation}
          Je mehr Nachrichten unter ständigen Störungen erfolgreich verarbeitet werden, desto sicherer ist nachgewiesen, dass Fehler im regulären Betrieb nahezu ausgeschlossen sind.
\end{enumerate}

\chapter{N}

Um die Einhaltung der Delivery-Garantien statistisch fundiert nachzuweisen, muss eine hinreichend große Anzahl an Alarmmeldungen verarbeitet werden. Da stochastische Tests das Auftreten von Duplikaten ($N_{\text{dup}}$) und Sequenzfehlern ($N_{\text{ooo}}$) prinzipbedingt nicht absolut ausschließen können, wird das statistische Verfahren der \textit{Rule of Three} nach Hanley und Lippman-Hand angewendet.

Beobachtet man in einer Testreihe von $N$ unabhängigen Übertragungen null Fehlereignisse ($N_{\text{dup}} = 0$ bzw. $N_{\text{ooo}} = 0$), so lässt sich mit einem Konfidenzniveau von 95\,\% eine obere Schranke für die wahre Fehlerrate $p$ pro Alarmmeldung ableiten:
\begin{equation}
    p \le \frac{3}{N} \quad \Longleftrightarrow \quad N \ge \frac{3}{p}
\end{equation}

Für die Dimensionierung des Testumfangs $N_{\text{msg}}$ dient der europäische Sicherheitsstandard \mbox{EN IEC 61784-3} (bzw. \mbox{EN 61508}) als Referenzrahmen. Für sicherheitskritische Alarmübertragungsanlagen im Brand- und Personenschutz (\mbox{SIL 2}) fordert die Norm eine maximale Restfehlerwahrscheinlichkeit von $p \le 10^{-8}$ pro Einzelnachricht für Übertragungsfehler wie Wiederholungen und Vertauschungen. Daraus ergibt sich für die experimentelle Evaluation eine geforderte Stichprobengröße von:
\begin{equation}
    N_{\text{msg}} \ge \frac{3}{10^{-8}} = 3 \cdot 10^{8} = 300\,000\,000 \text{ Alarmmeldungen}
\end{equation}

Das Kriterium für eine erfolgreiche Validierung der Implementierung ist somit erreicht, wenn über den Verlauf von $3 \cdot 10^8$ erfolgreich unter Fehlereinfluss übertragenen Alarmmeldungen gilt:
\begin{equation}
    N_{\text{dup}} = 0 \quad \text{und} \quad N_{\text{ooo}} = 0
\end{equation}
Damit ist mit einer Konfidenz von 95\,\% belegt, dass die Fehlerrate von \gls{alps} die geforderte Schranke von $10^{-8}$ nicht überschreitet.

\chapter{M}
Um sicher zu sein, dass die Garantien wirklich eingehalten werden, reicht es nicht, nur ein paar Alarme zu testen. Da Tests mit Zufallsfehlern niemals zu 100\,\% beweisen können, dass Fehler wie doppelte ($N_{\text{dup}}$) oder vertauschte Alarme ($N_{\text{ooo}}$) unmöglich sind, wird die sogenannte \textit{Rule of Three} verwendet.

Diese Regel besagt: Wenn in einem Test mit $N$ Nachrichten genau null Fehler auftreten, kann man mit 95\,\% Sicherheit sagen, dass das verbleibende Fehlerrisiko $p$ pro Alarm nicht größer ist als:
\begin{equation}
    p \le \frac{3}{N} \quad \text{bzw.} \quad N \ge \frac{3}{p}
\end{equation}

Als Vorgabe für den maximal erlaubten Fehlerwert $p$ dient die europäische Sicherheitsnorm \mbox{EN IEC 61784-3}. Für Alarmübertragungen im Bereich Brand- und Personenschutz (Sicherheitsstufe \mbox{SIL 2}) darf das Risiko für Duplikate oder Vertauschungen höchstens bei $p = 10^{-8}$ pro Nachricht liegen (also höchstens ein Fehler auf 100 Millionen Alarme). Setzt man diesen Wert in die Formel ein, ergibt sich die Anzahl der Alarme, die fehlerfrei übertragen werden müssen:
\begin{equation}
    N_{\text{msg}} \ge \frac{3}{10^{-8}} = 300\,000\,000 \text{ Alarmmeldungen}
\end{equation}

Der Test gilt als bestanden, wenn während der Übertragung dieser $300\,000\,000$ Alarme unter Störungseinfluss Folgendes gilt:
\begin{equation}
    N_{\text{dup}} = 0 \quad \text{und} \quad N_{\text{ooo}} = 0
\end{equation}
Damit ist mit 95\,\% statistischer Sicherheit nachgewiesen, dass \gls{alps} die Anforderungen von \mbox{SIL 2} erfüllt.